% ── methods/pretrain.tex ──────────────────────────────────────────────────────
\subsection{JEPA Pre-Training}
\label{sec:pretrain}

\paragraph{Architecture.}
The context encoder $f_\theta$ is a Transformer~\citep{vaswani} with
$d{=}384$, $8$ attention heads, $8$ layers, and $d_\text{ff}{=}1536$,
operating on amino acid sequences tokenised with a 25-token vocabulary
(20 canonical AAs + BOS / EOS / PAD / MASK / UNK).
An EMA target encoder $\bar{f}_\theta$ (momentum $= 0.996$) produces stable
representations without collapse.
A two-layer Transformer predictor $g_\phi$ bridges context and target
representations.

\paragraph{Objective.}
Given a sequence $\mathbf{x}$, we partition tokens into a context set
$\mathcal{C}$ (observed prefix) and a target set $\mathcal{T}$
(held-out suffix, 30--70\% of tokens).
$g_\phi$ maps context representations to target representations in
\emph{embedding space}:
\begin{equation}
  \mathcal{L}_\text{JEPA} =
  \mathbb{E}\bigl[\|g_\phi(f_\theta(\mathbf{x}_\mathcal{C})) -
  \bar{f}_\theta(\mathbf{x}_\mathcal{T})\|_2^2\bigr].
  \label{eq:jepa}
\end{equation}
Unlike masked language models, the optimisation target is a representation
rather than a discrete token identity.
This objective predicts in embedding space rather than token space.
We later test whether this leads to different fine-tuning behaviour
(Section~\ref{sec:res_emb}).

\paragraph{Training corpus.}
We pre-train on \textbf{868{,}724} unique AMP sequences ($5$--$50$ residues;
mean $35.4$~AA) drawn from APD3, DRAMP, UniProt AMP annotations (reviewed +
TrEMBL), and the AMPSphere experimental database.
AMPlify test sequences are excluded to prevent data leakage.
Training: AdamW ($\eta = 10^{-3}$, cosine annealing), batch size 2048,
100 epochs, converging to val loss $= 0.244$ in embedding MSE.
